\section{Recomendación práctica: primero estabilizar el entorno, después hablar de eficiencia de entrenamiento}
Para quien empieza, el mayor valor de Docker no es «ser más avanzado», sino hacer que el entorno de entrenamiento sea reproducible. Sobre todo en un flujo de trabajo de GPU remota con IDE local, basta con definir con claridad la imagen, los montajes y el directorio de trabajo para eliminar de antemano la mayoría de los problemas de entorno.

\begin{knowledgebox}{Flujo de trabajo mínimo recomendado}
\begin{itemize}
  \item Primero usar la imagen oficial para verificar que las dependencias estén completas
  \item Después usar un directorio montado para guardar el código y las salidas de los experimentos
  \item Por último conectar VSCode Remote para el desarrollo cotidiano
\end{itemize}
\end{knowledgebox}

\subsection{Resumen de la sección}
Docker no es una herramienta accesoria, sino infraestructura de la ingeniería de aprendizaje automático. Resolver primero el problema de la reproducibilidad del entorno y después optimizar la velocidad de entrenamiento suele ser más conveniente.

\newpage
